\documentclass{article}
\usepackage[T1]{fontenc}

\usepackage[a4paper, margin=2.5cm]{geometry}
\usepackage[polish, bidi=basic, provide=*]{babel}
\babelprovide[import, onchar=ids fonts]{polish}

\usepackage{amsmath}
\usepackage{booktabs}
\usepackage{hyperref}

\title{Raport z badania porównawczego systemów rozpoznawania mowy}
\author{Adam Krzywicki - 272523 Andrzej Manderla - XXXXXX Filip Zmyślony - XXXXXX}
\date{\today}

\begin{document}

\maketitle

\section{Wstęp i cel badania}
Celem niniejszego badania jest ocena przydatności dwóch wybranych systemów rozpoznawania mowy -- Whisper AI oraz Vosk -- do obsługi projektowanego interfejsu konwersacyjnego. 

Systemy zostały ocenione pod kątem jakości rozpoznawania mowy, poprawności transkrypcji oraz możliwości technicznych ich integracji z docelowym rozwiązaniem.

\section{Metodologia i zbiór testowy}
\subsection{Zbiór danych}
Badanie przeprowadzono na wspólnym zbiorze testowym, składającym się z 200 wypowiedzi pochodzących od 40 różnych mówców. Zbiór został podzielony na następujące kategorie tematyczne:
\begin{itemize}
    \item Polecenia i pytania typowe dla interfejsu,
    \item Wypowiedzi spontaniczne,
    \item Wypowiedzi zawierające nazwy własne, liczby, daty i terminy specjalistyczne,
    \item Wypowiedzi w trudniejszych warunkach akustycznych.
\end{itemize}
Dla każdego nagrania przygotowano transkrypcję referencyjną.

\subsection{Kryteria oceny i metryki}
Zgodnie z przyjętą metodyką, systemy oceniono na podstawie następujących wskaźników:

\textbf{1. Poprawność transkrypcji (WER -- Word Error Rate):}
Wskaźnik błędu słów obliczono ze wzoru:
\begin{equation}
WER = \frac{S + D + I}{N} \times 100\%
\end{equation}
gdzie: $S$ -- liczba substytucji, $D$ -- liczba pominięć, $I$ -- liczba wstawień, $N$ -- całkowita liczba słów w transkrypcji referencyjnej.

\textbf{2. Poprawność rozpoznawania ważnych informacji:}
Dla nazw własnych, dat, liczb oraz określania intencji zastosowano wskaźnik:
\begin{equation}
AI = \frac{N_{\text{poprawne}}}{N_{\text{wszystkie}}} \times 100\%
\end{equation}

\textbf{3. Szybkość działania (RTF -- Real-Time Factor):}
Stosunek czasu przetwarzania do czasu trwania nagrania:
\begin{equation}
RTF = \frac{T_p}{T_n}
\end{equation}
gdzie: $T_p$ -- czas przetwarzania nagrania przez system ASR, $T_n$ -- czas trwania nagrania dźwiękowego.

\section{Wyniki badań}

\subsection{Ogólna jakość transkrypcji i szybkość działania}

Tabela \ref{tab:ogolne} przedstawia wyniki dla 189 nagrań. Łączny czas nagrań wyniósł 2422,56~s,
a transkrypcje referencyjne zawierały 4313 słów.

\begin{table}[htbp]
\centering
\caption{Globalne wyniki WER i RTF dla badanych konfiguracji}
\label{tab:ogolne}
\begin{tabular}{lrrrrrr}
\toprule
\textbf{Konfiguracja} & \textbf{WER [\%]} & \textbf{S} & \textbf{D} &
\textbf{I} & \textbf{RTF} & \textbf{Czas [s]} \\
\midrule
Whisper \texttt{base}     & 42,50 & 1338 & 292 & 203 & 0,0503 & 121,81 \\
Whisper \texttt{large-v3} & 15,70 &  420 & 205 &  52 & 0,1415 & 342,68 \\
Vosk                     & 33,57 &  935 & 265 & 248 & 0,0792 & 191,75 \\
\bottomrule
\end{tabular}
\end{table}

Whisper \texttt{large-v3} osiągnął najniższy WER. W porównaniu
z wariantem \texttt{base} obniżył go o 26,80 punktu procentowego,
a przetworzenie zbioru zajęło około 2,8 raza więcej czasu.
Wszystkie konfiguracje uzyskały globalny RTF poniżej~1.

\subsection{Wyniki w poszczególnych kategoriach}

Tabela \ref{tab:kategorie} przedstawia średni WER w poszczególnych
kategoriach wypowiedzi. Dla każdej kategorii obliczono
średnią arytmetyczną wartości WER poszczególnych nagrań.

\begin{table}[htbp]
\centering
\caption{Średni WER według kategorii wypowiedzi [\%]}
\label{tab:kategorie}
\begin{tabular}{lrrr}
\toprule
\textbf{Kategoria wypowiedzi} &
\textbf{Whisper \texttt{base}} &
\textbf{Whisper \texttt{large-v3}} &
\textbf{Vosk} \\
\midrule
Polecenia i pytania (1)         & 33,43 & \textbf{6,53}  & 20,06 \\
Polecenia i pytania (2)         & 30,65 & \textbf{10,52} & 25,83 \\
Wypowiedzi spontaniczne         & 37,05 & \textbf{16,42} & 32,94 \\
Nazwy własne, liczby i daty     & 45,31 & \textbf{16,47} & 39,89 \\
Hałas w tle                     & 68,10 & \textbf{18,85} & 49,08 \\
\bottomrule
\end{tabular}
\end{table}

Whisper \texttt{large-v3} osiągnął najniższy średni WER we wszystkich
kategoriach. Największa różnica między konfiguracjami wystąpiła
w kategorii nagrań z hałasem w tle. W tej grupie WER modelu
\texttt{large-v3} był niższy o 49,25 punktu procentowego od wyniku
modelu \texttt{base} i o 30,23 punktu procentowego od wyniku Voska.

\subsection{Ocena wyodrębniania intencji i informacji kluczowych}

Rozpoznawanie intencji oceniono automatyczną heurystyką słów kluczowych
(\texttt{classify\_intent} w \texttt{bench.py}), bo tematyka nagrań jest zbyt
zróżnicowana (elektronika, gastronomia, hotel, zdrowie, motoryzacja...), żeby
zdefiniować jedną listę intencji domenowych. Wypowiedź klasyfikowana jest jako
jeden z typów aktu mowy: \emph{zamówienie}, \emph{rezerwacja}, \emph{problem},
\emph{polecenie}, \emph{pytanie} albo \emph{inne}, w tej kolejności priorytetu.
Intencję referencji i intencję hipotezy porównuje się, a zgodność liczona jest
jako \% nagrań, dla których obie intencje się zgadzają (kolumna
\texttt{Zgodność intencji [\%]} w \texttt{podsumowanie\_asr-*.csv}).

To jest heurystyka, nie model NLU ani ręczna adnotacja -- traktować jako
przybliżony wskaźnik, nie rozstrzygającą ocenę.

\begin{table}[htbp]
\centering
\caption{Zgodność intencji i dokładność rozpoznawania informacji kluczowych [\%]}
\label{tab:intencje}
\begin{tabular}{lccc}
\toprule
\textbf{System} & \textbf{Zgodność intencji} & \textbf{Nazwy własne} & \textbf{Daty/czas/liczby} \\
\midrule
Whisper \texttt{base}     & [UZUPEŁNIĆ] & [UZUPEŁNIĆ] & [UZUPEŁNIĆ] \\
Whisper \texttt{large-v3} & [UZUPEŁNIĆ] & [UZUPEŁNIĆ] & [UZUPEŁNIĆ] \\
Vosk                     & [UZUPEŁNIĆ] & [UZUPEŁNIĆ] & [UZUPEŁNIĆ] \\
\bottomrule
\end{tabular}
\end{table}

[Miejsce na analizę jakościową: przykłady błędów intencji i rozpoznawania
nazw/liczb po pełnym przebiegu \texttt{bench.py} -- dane z
\texttt{podsumowanie\_asr-*.csv}, kolumny \texttt{Zgodność intencji [\%]},
\texttt{Dokładność nazw własnych [\%]}, \texttt{Dokładność encji
(liczby/daty/czas) [\%]}.]

Uwaga: Tabela~\ref{tab:intencje} to automatyczna ocena na \emph{całym} zbiorze
(heurystyka), licząca inaczej niż ręczna ocena 36 nagrań w
Tabeli~\ref{tab:informacje-kluczowe} (ta druga pozostaje jako niezależna
weryfikacja na mniejszej, ręcznie sprawdzonej próbce).

\subsection{Poprawność rozpoznawania informacji kluczowych}

Przeprowadzono ręczną ocenę transkrypcji 36 nagrań z kategorii
\emph{Nazwy własne, liczby, daty}. W referencjach oznaczono
60 wystąpień nazw własnych i 99 wystąpień dat, określeń czasu lub liczb.
Wielowyrazową nazwę oceniano jako całość. Równoważne zapisy tej samej
wartości, np.\ \emph{24 października} i \emph{dwudziestego czwartego
października}, uznawano za poprawne.

\begin{table}[htbp]
\centering
\caption{Poprawność rozpoznawania informacji kluczowych}
\label{tab:informacje-kluczowe}
\begin{tabular}{lcc}
\toprule
\textbf{System} & \textbf{Nazwy własne} &
\textbf{Daty, czas i liczby} \\
\midrule
Whisper \texttt{base}     & 20/60 (33,33\%) & 69/99 (69,70\%) \\
Whisper \texttt{large-v3} & 47/60 (78,33\%) & 98/99 (98,99\%) \\
Vosk                     & 31/60 (51,67\%) & 81/99 (81,82\%) \\
\bottomrule
\end{tabular}
\end{table}

Whisper \texttt{large-v3} uzyskał najwyższą poprawność w obu
kryteriach. Przykładem jego błędu było rozpoznanie zakresu
ogniskowych obiektywu \emph{18--45 mm} jako \emph{18--35 mm}.

\subsection{Możliwości integracji}
Oba systemy działają lokalnie i mają gotowe biblioteki w Pythonie, co gwarantuje łatwą integrację oraz brak opłat za korzystanie z API. Dodatkowym bonusem jest pełna prywatność danych. Najważniejsza różnica leży w sprzęcie: Whisper potrzebuje większej mocy obliczeniowej, więc potrzebna jest w miarę mocna karta graficzna, podczas gdy Vosk jest lekki i bez problemu radzi sobie na zwykłym procesorze. Dodatkowo Vosk wspiera natywne przetwarzanie strumieniowe (audio podawane kawałkami w czasie rzeczywistym), podczas gdy Whisper w tym projekcie przetwarza cały plik na raz (batch) -- ma to konsekwencje dla opóźnienia, patrz sekcja~\ref{sec:latency}. Pełne porównanie API, formatu wyników, wymagań sprzętowych i kosztów -- \texttt{docs/INTEGRACJA.md}.

\subsection{Porównanie sparowane per nagranie}

Dodatkowo do porównania na poziomie zagregowanych statystyk, dla każdego
nagrania policzono różnicę WER między systemami (\texttt{scripts/paired\_comparison.py},
wynik w \texttt{wyniki\_asr-*\_parowane\_*.csv}) -- to pozwala stwierdzić, na
ilu konkretnych nagraniach jeden system wypada lepiej od drugiego, a nie tylko
jaka jest średnia różnica w całej próbce.

\begin{table}[htbp]
\centering
\caption{Porównanie sparowane Whisper vs Vosk (per nagranie)}
\label{tab:parowane}
\begin{tabular}{lccc}
\toprule
\textbf{Para systemów} & \textbf{Vosk lepszy} & \textbf{Vosk gorszy} & \textbf{Remis} \\
\midrule
Whisper \texttt{base} vs Vosk     & [UZUPEŁNIĆ] & [UZUPEŁNIĆ] & [UZUPEŁNIĆ] \\
Whisper \texttt{large-v3} vs Vosk & [UZUPEŁNIĆ] & [UZUPEŁNIĆ] & [UZUPEŁNIĆ] \\
\bottomrule
\end{tabular}
\end{table}

[Miejsce na komentarz: średnia/mediana różnicy WER (B - A) i wynik testu
Wilcoxona (jeśli dostępny scipy), z \texttt{scripts/paired\_comparison.py}.]

\subsection{Pomiar opóźnienia w rzeczywistym przepływie}
\label{sec:latency}

Globalny RTF (sekcja~3.3) mówi, jak szybko system przetwarza nagranie
\emph{względem jego długości}, ale nie odpowiada na praktyczne pytanie: ile
użytkownik czeka \emph{po tym, jak przestanie mówić}, zanim dostanie
transkrypcję. Zmierzono to na reprezentatywnym podzbiorze (po kilka nagrań z
każdej kategorii, \texttt{scripts/latency\_pipeline.py}), symulując
rzeczywisty przepływ: audio "dociera" w czasie rzeczywistym (z opóźnieniem
odpowiadającym długości nagrania), a mierzone jest tylko opóźnienie
\emph{ogonowe} -- czas od ostatniego fragmentu audio do gotowej transkrypcji.

Vosk przetwarza audio strumieniowo (kawałkami, w miarę jak "nadchodzą"), więc
jego opóźnienie ogonowe jest małe i niezależne od długości nagrania. Whisper
w tej konfiguracji nie ma API przyrostowego -- może zacząć transkrybować
dopiero, gdy całe nagranie jest już dostępne, więc jego opóźnienie ogonowe
równa się czasowi pojedynczego wywołania \texttt{transcribe()}.

\begin{table}[htbp]
\centering
\caption{Opóźnienie ogonowe (od końca wypowiedzi do transkrypcji)}
\label{tab:latency}
\begin{tabular}{lcc}
\toprule
\textbf{System} & \textbf{Mediana [s]} & \textbf{Maksimum [s]} \\
\midrule
Whisper \texttt{base}     & [UZUPEŁNIĆ] & [UZUPEŁNIĆ] \\
Whisper \texttt{large-v3} & [UZUPEŁNIĆ] & [UZUPEŁNIĆ] \\
Vosk                     & [UZUPEŁNIĆ] & [UZUPEŁNIĆ] \\
\bottomrule
\end{tabular}
\end{table}

[Dane z \texttt{results/latency\_pipeline.csv} po pełnym przebiegu
\texttt{scripts/latency\_pipeline.py}.]

\subsection{Ocena końcowa (synteza ważona)}

Wyniki wszystkich 9 kryteriów przeskalowano do wspólnej skali 0--100 (wyżej =
lepiej) i połączono w jedną ocenę ważoną wagami z wytycznych badania
(\texttt{scripts/compute\_final\_score.py}, metodologia przeskalowania opisana
w docstringu skryptu):

\begin{itemize}
    \item WER, WER podzbioru spontaniczne+pozawerbalne: $100 - WER$,
    \item Nazwy własne, daty/czas/liczby, intencja, kompletność: wartość \% już 0--100,
    \item Odporność na szum: $100 - \Delta WER_{szum}$,
    \item Szybkość (RTF): $100 \times (1 - RTF)$, $RTF{=}1$ (czas rzeczywisty) $\to 0$,
    \item Integracja: ocena opisowa (0--100) na podstawie \texttt{docs/INTEGRACJA.md}, podawana ręcznie -- nie da się jej wyliczyć automatycznie.
\end{itemize}

\begin{table}[htbp]
\centering
\caption{Ocena końcowa (ważona synteza 9 kryteriów, skala 0--100)}
\label{tab:ocena-koncowa}
\begin{tabular}{lc}
\toprule
\textbf{System} & \textbf{Ocena końcowa} \\
\midrule
Whisper \texttt{base}     & [UZUPEŁNIĆ] \\
Whisper \texttt{large-v3} & [UZUPEŁNIĆ] \\
Vosk                     & [UZUPEŁNIĆ] \\
\bottomrule
\end{tabular}
\end{table}

[Pełna tabela z rozbiciem na poszczególne kryteria: \texttt{results/podsumowanie\_asr-*\_ocena\_koncowa.csv}.]

\section{Podsumowanie i wnioski końcowe}

Na podstawie przeprowadzonych pomiarów i analiz sformułowano następujące wnioski:

\begin{enumerate}
    \item \textbf{Rozpoznawanie wypowiedzi charakterystycznych dla interfejsu:} \\
    Zdecydowanie najdokładniejszym systemem jest Whisper w konfiguracji large-v3. Osiągnął on błąd na poziomie zaledwie 6,53\% -- 10,52\% w kategoriach typowych poleceń i pytań, co pozwala na niezawodne dekodowanie intencji użytkownika. Vosk z wynikiem 20\%--25\% również stanowi akceptowalne rozwiązanie, choć częściej może prowadzić do nieporozumień.
    
    \item \textbf{Różnice w kategoriach i warunkach akustycznych:} \\
    Najbardziej drastyczne różnice zauważono w nagraniach z hałasem w tle oraz przy nazwach własnych. Whisper large-v3 wykazał niezwykłą odporność na szum (WER na poziomie 18,85\%), podczas gdy Vosk (49,08\%) i Whisper base (68,10\%) ulegały znacznej degradacji. 
    
    \item \textbf{Relacja jakości do czasu przetwarzania:} \\
    Zwiększenie dokładności transkrypcji w modelu large-v3 kosztuje -- system przetwarzał dane około 2,8 raza dłużej niż bazowa wersja Whispera. Mimo to, wskaźnik RTF na poziomie 0,1415 oznacza, że wygenerowanie tekstu następuje wielokrotnie szybciej niż trwa sama wypowiedź. Takie opóźnienie jest w pełni akceptowalne i niezauważalne dla użytkownika końcowego projektowanego interfejsu.
    
    \item \textbf{Rekomendacja wdrożeniowa:} \\
    Ostateczny wybór zależy od dostępnej architektury sprzętowej:
    \begin{itemize}
        \item Jeśli dysponujemy serwerem obliczeniowym wyposażonym w odpowiednio silną kartę graficzną, bezwzględnie rekomenduje się wdrożenie systemu Whisper large-v3. Gwarantuje on najwyższą skuteczność.
        \item Jeśli zależy nam na zminimalizowaniu kosztów sprzętowych i planujemy uruchamiać interfejs bezpośrednio na urządzeniach wbudowanych bez dostępu do GPU, rekomendowanym wyborem jest Vosk. Choć jest mniej precyzyjny, jego działanie na samym procesorze przy jednoczesnym braku kosztów chmurowych i zachowaniu całkowitej prywatności stanowi świetny kompromis dla mniej wymagających zastosowań.
    \end{itemize}
\end{enumerate}

\end{document}